\chapter{1994年旧科目组卷（理）}

\section{单选题}

\begin{problem}
极坐标方程 \(\rho = \cos \left(\frac{\pi}{4} - \theta\right)\) 所表示的曲线是（\quad）

\choices
    {双曲线}
    {椭圆}
    {抛物线}
    {圆}
\end{problem}

\begin{answer}
D
\end{answer}

\begin{solution}
由 \(\cos \left(\frac{\pi}{4} - \theta\right) = \frac{1}{\sqrt{2}}\left(\cos \theta + \sin \theta\right)\)，得 \(\rho = \frac{1}{\sqrt{2}}\left(\cos \theta + \sin \theta\right)\). 两边同乘以 \(\rho\)，并利用 \(x = \rho \cos \theta\)，\(y = \rho \sin \theta\)，得 \(x^2 + y^2 = \frac{x + y}{\sqrt{2}}\). 配方得 \(\left(x - \frac{1}{2\sqrt{2}}\right)^2 + \left(y - \frac{1}{2\sqrt{2}}\right)^2 = \frac{1}{4}\)，这是一个圆，故选 D.
\end{solution}

\begin{problem}
如果方程 \(x^2 + ky^2 = 2\) 表示焦点在 \(y\) 轴上的椭圆，那么实数 \(k\) 的取值范围是（\quad）

\choices
    {\((0, +\infty)\)}
    {\((0, 2)\)}
    {\((1, +\infty)\)}
    {\((0, 1)\)}
\end{problem}

\begin{answer}
D
\end{answer}

\begin{solution}
方程表示椭圆，必须有 \(k > 0\)，此时方程化为 \(\frac{x^2}{2} + \frac{y^2}{\frac{2}{k}} = 1\). 焦点在 \(y\) 轴上，说明 \(y\) 方向的半轴更长，即 \(\frac{2}{k} > 2\)，解得 \(0 < k < 1\)，故选 D.
\end{solution}

\begin{problem}
\(\displaystyle\lim_{n \to \infty} \frac{\mathrm C_n^0 + \mathrm C_n^1 + \mathrm C_n^2 + \cdots + \mathrm C_n^n}{1 + 2 + 2^2 + \cdots + 2^n} =\)（\quad）

\choices
    {\(0\)}
    {\(\frac{1}{2}\)}
    {\(1\)}
    {\(2\)}
\end{problem}

\begin{answer}
B
\end{answer}

\begin{solution}
分子是二项式系数之和 \(\mathrm C_n^0 + \mathrm C_n^1 + \cdots + \mathrm C_n^n = 2^n\)，分母是等比数列之和 \(1 + 2 + \cdots + 2^n = 2^{n+1} - 1\). 因此所求极限为 \(\displaystyle\lim_{n \to \infty} \frac{2^n}{2^{n+1} - 1} = \frac{1}{2}\)，故选 B.
\end{solution}

\begin{problem}
设 \(\theta\) 是第二象限的角，则必有（\quad）

\choices
    {\(\tan \frac{\theta}{2} > \cot \frac{\theta}{2}\)}
    {\(\tan \frac{\theta}{2} < \cot \frac{\theta}{2}\)}
    {\(\sin \frac{\theta}{2} > \cos \frac{\theta}{2}\)}
    {\(\sin \frac{\theta}{2} < \cos \frac{\theta}{2}\)}
\end{problem}

\begin{answer}
A
\end{answer}

\begin{solution}
因为 \(\theta\) 是第二象限的角，所以存在 \(k \in \Z\)，使 \(\frac{\pi}{2} + 2k\pi < \theta < \pi + 2k\pi\). 令 \(\varphi = \frac{\theta}{2}\)，则 \(\frac{\pi}{4} + k\pi < \varphi < \frac{\pi}{2} + k\pi\). 由正切函数以 \(\pi\) 为周期，得 \(\tan \varphi > 1\)，从而 \(\cot \varphi = \frac{1}{\tan \varphi} < 1\)，所以 \(\tan \frac{\theta}{2} > \cot \frac{\theta}{2}\) 恒成立.
另一方面，\(\sin \varphi - \cos \varphi = \sqrt{2}\sin \left(\varphi - \frac{\pi}{4}\right)\)，而 \(\varphi - \frac{\pi}{4} \in \left(k\pi, k\pi + \frac{\pi}{4}\right)\) 的正负随 \(k\) 的奇偶而变，故 C、D 均不恒成立，故选 A.
\end{solution}

\begin{problem}
若直线 \(x + ay + 2 = 0\) 和 \(2x + 3y + 1 = 0\) 互相垂直，则 \(a =\)（\quad）

\choices
    {\(-\frac{2}{3}\)}
    {\(-\frac{3}{2}\)}
    {\(\frac{2}{3}\)}
    {\(\frac{3}{2}\)}
\end{problem}

\begin{answer}
A
\end{answer}

\begin{solution}
两直线的方向向量分别可取 \((a, -1)\) 与 \((3, -2)\)，互相垂直的条件是数量积为零，即 \(3a + 2 = 0\)，解得 \(a = -\frac{2}{3}\)，故选 A.
\end{solution}

\begin{problem}
某种细菌在培养过程中，每 \(20\) 分钟分裂一次（一个分裂为两个）. 经过 \(3\) 小时，这种细菌由 \(1\) 个可繁殖成（\quad）

\choices
    {\(511\) 个}
    {\(512\) 个}
    {\(1023\) 个}
    {\(1024\) 个}
\end{problem}

\begin{answer}
B
\end{answer}

\begin{solution}
\(3\) 小时等于 \(180\) 分钟，共分裂 \(\frac{180}{20} = 9\) 次. 每次分裂后数量变为原来的 \(2\) 倍，所以最后有 \(1 \cdot 2^9 = 512\) 个，故选 B.
\end{solution}

\begin{problem}
在下列函数中，以 \(\frac{\pi}{2}\) 为周期的函数是（\quad）

\choices
    {\(y = \sin 2x + \cos 4x\)}
    {\(y = \sin 2x \cos 4x\)}
    {\(y = \sin 2x + \cos 2x\)}
    {\(y = \sin 2x \cos 2x\)}
\end{problem}

\begin{answer}
D
\end{answer}

\begin{solution}
对 D，有 \(\sin 2x \cos 2x = \frac{1}{2}\sin 4x\)，而 \(\sin 4x\) 的最小正周期为 \(\frac{2\pi}{4} = \frac{\pi}{2}\)，所以 D 以 \(\frac{\pi}{2}\) 为周期.
A 中的 \(\sin 2x\) 以 \(\pi\) 为最小正周期：若 A 以 \(\frac{\pi}{2}\) 为周期，则 \(\sin 2\left(x + \frac{\pi}{2}\right) + \cos 4\left(x + \frac{\pi}{2}\right) = \sin 2x + \cos 4x\) 对一切 \(x\) 成立，即 \(-\sin 2x = \sin 2x\)，不可能. 同理 B 中的因子 \(\sin 2x\) 与 C 中的 \(\sin 2x + \cos 2x\) 的最小正周期都是 \(\pi\)，都不以 \(\frac{\pi}{2}\) 为周期，故选 D.
\end{solution}

\begin{problem}
已知正六棱台的上、下底面边长分别为 \(2\) 和 \(4\)，高为 \(2\)，则其体积为（\quad）

\choices
    {\(32\sqrt{3}\)}
    {\(28\sqrt{3}\)}
    {\(24\sqrt{3}\)}
    {\(20\sqrt{3}\)}
\end{problem}

\begin{answer}
B
\end{answer}

\begin{solution}
边长为 \(a\) 的正六边形面积为 \(6 \cdot \frac{\sqrt{3}}{4}a^2 = \frac{3\sqrt{3}}{2}a^2\). 于是上底面积 \(S_1 = 6\sqrt{3}\)，下底面积 \(S_2 = 24\sqrt{3}\). 由棱台体积公式 \(V = \frac{h}{3}\left(S_1 + S_2 + \sqrt{S_1 S_2}\right)\)，得 \(V = \frac{2}{3}\left(6\sqrt{3} + 24\sqrt{3} + 12\sqrt{3}\right) = 28\sqrt{3}\)，故选 B.
\end{solution}

\begin{problem}
使 \(\left(\sqrt{6} - \sqrt{2}\i\right)^n\) 是纯虚数的最小自然数 \(n =\)（\quad）

\choices
    {\(3\)}
    {\(4\)}
    {\(5\)}
    {\(6\)}
\end{problem}

\begin{answer}
A
\end{answer}

\begin{solution}
记 \(z = \sqrt{6} - \sqrt{2}\i\)，则 \(|z| = \sqrt{6 + 2} = 2\sqrt{2}\)，且 \(\arg z = -\frac{\pi}{6}\). 于是 \(z^n\) 的辐角为 \(-\frac{n\pi}{6}\)（相差 \(2\pi\) 整数倍）. \(z^n\) 为纯虚数当且仅当 \(-\frac{n\pi}{6} = \frac{\pi}{2} + k\pi\)（\(k \in \Z\)），即 \(n = -3 - 6k\). 使 \(n\) 为自然数的最小值是 \(n = 3\)（对应 \(k = -1\)），故选 A.
\end{solution}

\begin{problem}
有甲、乙、丙三项任务，甲需 \(2\) 人承担，乙、丙各需 \(1\) 人承担. 从 \(10\) 人中选派 \(4\) 人承担这三项任务，不同的选法共有（\quad）

\choices
    {\(1260\) 种}
    {\(2025\) 种}
    {\(2520\) 种}
    {\(5040\) 种}
\end{problem}

\begin{answer}
C
\end{answer}

\begin{solution}
先从 \(10\) 人中选出承担任务的 \(4\) 人，再从这 \(4\) 人中选出 \(2\) 人承担甲，最后将剩下的 \(2\) 人分别安排给乙、丙. 因此选法数为 \(\mathrm C_{10}^{4}\mathrm C_{4}^{2}\mathrm A_{2}^{2} = 210 \cdot 6 \cdot 2 = 2520\)，故选 C.
\end{solution}

\begin{problem}
对于直线 \(m\)，\(n\) 和平面 \(\alpha\)，\(\beta\)，\(\alpha \perp \beta\) 的一个充分条件是（\quad）

\choices
    {\(m \perp n\)，\(m \parallel \alpha\)，\(n \parallel \beta\)}
    {\(m \perp n\)，\(\alpha \cap \beta = m\)，\(n \subset \alpha\)}
    {\(m \parallel n\)，\(n \perp \beta\)，\(m \subset \alpha\)}
    {\(m \parallel n\)，\(m \perp \alpha\)，\(n \perp \beta\)}
\end{problem}

\begin{answer}
C
\end{answer}

\begin{solution}
对 C，由 \(n \perp \beta\) 且 \(m \parallel n\) 可知 \(m \perp \beta\). 又 \(m \subset \alpha\)，所以平面 \(\alpha\) 内有一条直线 \(m\) 垂直于平面 \(\beta\)，根据平面与平面垂直的判定定理得 \(\alpha \perp \beta\).
A 只给出直线之间的垂直与分别平行关系，不能推出两个平面垂直. B 中 \(n \subset \alpha\) 且 \(n \perp m\) 不能保证 \(n\) 垂直于平面 \(\beta\). D 中由 \(m \parallel n\)，\(m \perp \alpha\) 得 \(n \perp \alpha\)，再结合 \(n \perp \beta\) 得到的是 \(\alpha \parallel \beta\). 因此只有 C 是充分条件，故选 C.
\end{solution}

\begin{problem}
设函数 \(f(x) = 1 - \sqrt{1 - x^2}\)（\(-1 \leqslant x \leqslant 0\)），则函数 \(y = f^{-1}(x)\) 的图象是（\quad）

\choices
{\choicebitmap[width=0.15\paperwidth]{img/1994/old_subject_science/q12_optA.png}}
{\choicebitmap[width=0.15\paperwidth]{img/1994/old_subject_science/q12_optB.png}}
{\choicebitmap[width=0.15\paperwidth]{img/1994/old_subject_science/q12_optC.png}}
{\choicebitmap[width=0.15\paperwidth]{img/1994/old_subject_science/q12_optD.png}}
\end{problem}

\begin{answer}
B
\end{answer}

\begin{solution}
由 \(y = 1 - \sqrt{1 - x^2}\)（\(-1 \leqslant x \leqslant 0\)）得 \(\sqrt{1 - x^2} = 1 - y\)，故 \(0 \leqslant y \leqslant 1\)，且 \(f\) 在 \([-1, 0]\) 上严格递减. 反函数的图象与原函数图象关于直线 \(y = x\) 对称，所以反函数的定义域为 \([0, 1]\)，值域为 \([-1, 0]\)，图象是从原点出发到达点 \((1, -1)\) 的一段弧，排除 A、C.
将 \(x\) 与 \(y\) 互换，得反函数满足 \((x - 1)^2 + y^2 = 1\)（\(0 \leqslant x \leqslant 1\)，\(y \leqslant 0\)）. 在点 \((0, 0)\) 附近，\(y = -\sqrt{2x - x^2}\) 关于 \(x\) 的变化率趋于无穷，曲线以竖直切线离开原点，与 B 相符（D 在原点处的切线是水平的），故选 B.
\end{solution}

\begin{problem}
已知过球面上 \(A\)，\(B\)，\(C\) 三点的截面和球心的距离等于球半径的一半，且 \(AB = BC = CA = 2\)，则球面面积是（\quad）

\choices
    {\(\frac{16}{9}\pi\)}
    {\(\frac{8}{3}\pi\)}
    {\(4\pi\)}
    {\(\frac{64}{9}\pi\)}
\end{problem}

\begin{answer}
D
\end{answer}

\begin{solution}
球心在截面上的射影是截面圆的圆心. 因为 \(ABC\) 是边长为 \(2\) 的正三角形，截面圆就是其外接圆，半径为 \(r = \frac{2}{\sqrt{3}}\). 设球半径为 \(R\)，球心到截面的距离为 \(\frac{R}{2}\)，则由截面性质 \(r^2 = R^2 - \left(\frac{R}{2}\right)^2 = \frac{3R^2}{4}\). 于是 \(\frac{4}{3} = \frac{3R^2}{4}\)，解得 \(R^2 = \frac{16}{9}\). 故球面面积为 \(4\pi R^2 = \frac{64}{9}\pi\)，故选 D.
\end{solution}

\begin{problem}
函数 \(y = \arccos \left(\sin x\right)\)（\(-\frac{\pi}{3} < x < \frac{2\pi}{3}\)）的值域是（\quad）

\choices
    {\(\left(\frac{\pi}{6}, \frac{5\pi}{6}\right)\)}
    {\(\left[0, \frac{5\pi}{6}\right)\)}
    {\(\left(\frac{\pi}{3}, \frac{2\pi}{3}\right)\)}
    {\(\left[\frac{\pi}{6}, \frac{2\pi}{3}\right)\)}
\end{problem}

\begin{answer}
B
\end{answer}

\begin{solution}
当 \(-\frac{\pi}{3} < x < \frac{2\pi}{3}\) 时，\(\sin x\) 的取值范围是 \(\left(-\frac{\sqrt{3}}{2}, 1\right]\)：最大值 \(1\) 在 \(x = \frac{\pi}{2}\) 处取到，左端点处 \(\sin \left(-\frac{\pi}{3}\right) = -\frac{\sqrt{3}}{2}\) 取不到. 因为 \(\arccos t\) 在 \([-1, 1]\) 上严格递减，所以 \(y\) 的最小值是 \(\arccos 1 = 0\)（可取到），且 \(y < \arccos \left(-\frac{\sqrt{3}}{2}\right) = \frac{5\pi}{6}\)（取不到）. 故值域为 \(\left[0, \frac{5\pi}{6}\right)\)，故选 B.
\end{solution}

\begin{problem}
定义在 \((- \infty, + \infty)\) 上的任意函数 \(f(x)\) 都可以表示成一个奇函数 \(g(x)\) 和一个偶函数 \(h(x)\) 之和. 如果 \(f(x) = \lg \left(10^x + 1\right)\)，\(x \in (-\infty, +\infty)\)，那么（\quad）

\choices
    {\(g(x) = x\)，\(h(x) = \lg \left(10^x + 10^{-x} + 2\right)\)}
    {\(g(x) = \frac{1}{2}\left[\lg \left(10^x + 1\right) + x\right]\)，\(h(x) = \frac{1}{2}\left[\lg \left(10^x + 1\right) - x\right]\)}
    {\(g(x) = \frac{x}{2}\)，\(h(x) = \lg \left(10^x + 1\right) - \frac{x}{2}\)}
    {\(g(x) = -\frac{x}{2}\)，\(h(x) = \lg \left(10^x + 1\right) + \frac{x}{2}\)}
\end{problem}

\begin{answer}
C
\end{answer}

\begin{solution}
记 \(f(x) = \lg \left(10^x + 1\right)\)，则 \(f(-x) = \lg \left(10^{-x} + 1\right) = \lg \frac{10^x + 1}{10^x} = f(x) - x\). 若 \(f = g + h\)，其中 \(g\) 为奇函数、\(h\) 为偶函数，则 \(f(-x) = -g(x) + h(x)\). 两式相加、相减得 \(g(x) = \frac{f(x) - f(-x)}{2} = \frac{x}{2}\)，\(h(x) = \frac{f(x) + f(-x)}{2} = f(x) - \frac{x}{2} = \lg \left(10^x + 1\right) - \frac{x}{2}\). 直接验证 \(\frac{x}{2}\) 为奇函数，而 \(h(-x) = f(-x) + \frac{x}{2} = f(x) - x + \frac{x}{2} = h(x)\) 为偶函数. 故选 C.
\end{solution}

\section{填空题}

\begin{problem}
抛物线 \(y^2 = 8 - 4x\) 的准线方程是\fillinblank{}.
\end{problem}

\begin{answer}
\(x = 3\)
\end{answer}

\begin{solution}
将方程写成 \(y^2 = -4(x - 2)\)，顶点为 \((2, 0)\)，开口向左，且 \(2p = 4\)，即 \(p = 2\). 准线是过点 \((2 + 1, 0)\) 且垂直于 \(x\) 轴的直线，故准线方程为 \(x = 3\).
\end{solution}

\begin{problem}
在 \((x + m)^7\)（\(m \in \N\)）的展开式中，\(x^5\) 的系数是 \(x^6\) 的系数与 \(x^4\) 的系数的等差中项，则 \(m =\fillinblank{}\).
\end{problem}

\begin{answer}
\(1\)
\end{answer}

\begin{solution}
\((x + m)^7\) 展开式的通项为 \(\mathrm C_7^k x^k m^{7-k}\). 于是 \(x^5\) 的系数为 \(21m^2\)，\(x^6\) 的系数为 \(7m\)，\(x^4\) 的系数为 \(35m^3\). 由等差中项条件 \(2 \cdot 21m^2 = 7m + 35m^3\)，即 \(5m^2 - 6m + 1 = 0\)（\(m > 0\)）. 解得 \(m = 1\) 或 \(m = \frac{1}{5}\)，结合 \(m \in \N\) 得 \(m = 1\).
\end{solution}

\begin{problem}
若 \(\frac{5}{4}\pi < \theta < \frac{3}{2}\pi\)，\(\sin 2\theta = a\)，则 \(\cos \theta - \sin \theta\) 的值是\fillinblank{}.
\end{problem}

\begin{answer}
\(\sqrt{1 - a}\)
\end{answer}

\begin{solution}
由 \(\left(\cos \theta - \sin \theta\right)^2 = 1 - \sin 2\theta = 1 - a\)，得 \(\cos \theta - \sin \theta = \pm\sqrt{1 - a}\). 当 \(\frac{5}{4}\pi < \theta < \frac{3}{2}\pi\) 时，例如取 \(\theta = \frac{4\pi}{3}\)，有 \(\cos \theta - \sin \theta = -\frac{1}{2} + \frac{\sqrt{3}}{2} > 0\). 一般地，此时 \(\theta - \frac{\pi}{4} \in (\pi, \frac{5\pi}{4})\)，而 \(\cos \theta - \sin \theta = \sqrt{2}\cos \left(\theta + \frac{\pi}{4}\right)\)，\(\theta + \frac{\pi}{4} \in \left(\frac{3\pi}{2}, \frac{7\pi}{4}\right)\) 的余弦为正. 故取正号，即 \(\cos \theta - \sin \theta = \sqrt{1 - a}\).
\end{solution}

\begin{problem}
设圆锥底面圆周上两点 \(A\)，\(B\) 间的距离为 \(2\)，圆锥顶点到直线 \(AB\) 的距离为 \(\sqrt{3}\)，\(AB\) 和圆锥的轴的距离为 \(1\)，则该圆锥的体积为\fillinblank{}.
\end{problem}

\begin{answer}
\(\frac{2\sqrt{2}}{3}\pi\)
\end{answer}

\begin{solution}
设圆锥底面圆心为 \(O\)，顶点为 \(S\)，\(M\) 为弦 \(AB\) 的中点. 由圆心与弦中点的连线垂直于弦，得 \(OM \perp AB\)，且 \(AM = 1\). 轴到直线 \(AB\) 的距离就是 \(OM = 1\)，于是在直角三角形 \(OAM\) 中，底面半径满足 \(r^2 = OA^2 = OM^2 + AM^2 = 2\).
又 \(SO\) 垂直于底面，故 \(SO \perp AB\). 结合 \(OM \perp AB\) 可知 \(AB\) 垂直于平面 \(SOM\)，从而 \(SM \perp AB\)，即 \(SM\) 就是顶点到直线 \(AB\) 的距离 \(\sqrt{3}\). 在直角三角形 \(SOM\) 中，圆锥的高满足 \(h^2 = SO^2 = SM^2 - OM^2 = 3 - 1 = 2\)，即 \(h = \sqrt{2}\). 因此圆锥体积为 \(\frac{1}{3}\pi r^2 h = \frac{1}{3}\pi \cdot 2 \cdot \sqrt{2} = \frac{2\sqrt{2}}{3}\pi\).
\end{solution}

\begin{problem}
在测量某物理量的过程中，因仪器和观察的误差，使得 \(n\) 次测量分别得到 \(a_1\)，\(a_2\)，\(\cdots\)，\(a_n\)，共 \(n\) 个数据. 我们规定所测量物理量的“最佳近似值”\(a\) 是这样一个量：与其他近似值比较，\(a\) 与各数据的差的平方和最小. 依此规定，从 \(a_1\)，\(a_2\)，\(\cdots\)，\(a_n\) 推出的 \(a =\fillinblank{}\).
\end{problem}

\begin{answer}
\(\frac{a_1 + a_2 + \cdots + a_n}{n}\)
\end{answer}

\begin{solution}
设近似值为 \(t\)，考察误差平方和 \(F(t) = \sum_{i=1}^{n}\left(t - a_i\right)^2\). 记 \(\bar{a} = \frac{a_1 + a_2 + \cdots + a_n}{n}\)，则 \(\sum_{i=1}^{n}\left(\bar{a} - a_i\right) = 0\). 于是 \(F(t) - F(\bar{a}) = \sum_{i=1}^{n}\left[\left(t - a_i\right)^2 - \left(\bar{a} - a_i\right)^2\right] = n\left(t - \bar{a}\right)^2 \geqslant 0\)，等号仅当 \(t = \bar{a}\) 时成立. 故最佳近似值为 \(a = \frac{a_1 + a_2 + \cdots + a_n}{n}\).
\end{solution}

\section{解答题}

\begin{problem}
已知 \(z = 1 + \i\).

\begin{enumerate}
\item 设 \(\omega = z^2 + 3\overline{z} - 4\)，求 \(\omega\) 的三角形式；
\item 如果 \(\frac{z^2 + az + b}{z^2 - z + 1} = 1 - \i\)，求实数 \(a\)，\(b\) 的值.
\end{enumerate}
\end{problem}

\begin{solution}
\begin{enumerate}
\item 由 \(z = 1 + \i\) 得 \(\overline{z} = 1 - \i\)，\(z^2 = 2\i\). 因此 \(\omega = z^2 + 3\overline{z} - 4 = 2\i + 3(1 - \i) - 4 = -1 - \i\). 于是 \(|\omega| = \sqrt{2}\)，且 \(\omega\) 位于第三象限，辐角可取 \(\frac{5\pi}{4}\). 故 \(\omega = \sqrt{2}\left(\cos \frac{5\pi}{4} + \i\sin \frac{5\pi}{4}\right)\).
\item 先计算分母 \(z^2 - z + 1 = 2\i - (1 + \i) + 1 = \i \neq 0\). 分子为 \(z^2 + az + b = 2\i + a(1 + \i) + b = (a + b) + (a + 2)\i\). 由题设，分子等于 \(\i(1 - \i) = 1 + \i\). 根据复数相等的定义，得 \(a + b = 1\)，\(a + 2 = 1\). 解得 \(a = -1\)，\(b = 2\).
\end{enumerate}
\end{solution}

\begin{problem}
已知函数 \(f(x) = \tan x\)，\(x \in \left(0, \frac{\pi}{2}\right)\). 若 \(x_1\)，\(x_2 \in \left(0, \frac{\pi}{2}\right)\)，且 \(x_1 \neq x_2\)，证明 \(\frac{1}{2}\left[f(x_1) + f(x_2)\right] > f\left(\frac{x_1 + x_2}{2}\right)\).
\end{problem}

\begin{solution}
由正切的和角公式，\(\tan x_1 + \tan x_2 = \frac{\sin x_1}{\cos x_1} + \frac{\sin x_2}{\cos x_2} = \frac{\sin \left(x_1 + x_2\right)}{\cos x_1 \cos x_2}\). 记 \(S = x_1 + x_2\)，\(D = x_1 - x_2\)，则 \(S \in (0, \pi)\)，\(D \neq 0\). 由积化和差，\(\cos x_1 \cos x_2 = \frac{1}{2}\left[\cos S + \cos D\right]\)，所以 \(\frac{1}{2}\left(\tan x_1 + \tan x_2\right) = \frac{\sin S}{\cos S + \cos D}\). 另一方面，\(\tan \frac{S}{2} = \frac{\sin S}{1 + \cos S}\). 两式相减得 \(\frac{1}{2}\left(\tan x_1 + \tan x_2\right) - \tan \frac{S}{2} = \frac{\sin S\left(1 - \cos D\right)}{\left(\cos S + \cos D\right)\left(1 + \cos S\right)}\). 其中 \(\sin S > 0\)（\(0 < S < \pi\)），\(1 - \cos D > 0\)（\(D \neq 0\)），\(\cos S + \cos D = 2\cos x_1 \cos x_2 > 0\)，\(1 + \cos S > 0\). 故上式为正，即 \(\frac{1}{2}\left[f(x_1) + f(x_2)\right] > f\left(\frac{x_1 + x_2}{2}\right)\).
\end{solution}

\begin{problem}
如图，已知 \(A_{1}B_{1}C_{1} - ABC\) 是正三棱柱，\(D\) 是 \(AC\) 中点.

\begin{enumerate}
\item 证明 \(AB_1 \parallel\) 平面 \(DBC_1\)；
\item 假设 \(AB_1 \perp BC_1\)，求以 \(BC_1\) 为棱，\(DBC_1\) 与 \(CBC_1\) 为面的二面角 \(\alpha\) 的度数.
\end{enumerate}

\bitmapfigure[width=0.5\linewidth]{img/1994/old_subject_science/q23_fig1.png}
\end{problem}

\begin{solution}
\begin{enumerate}
\item 因为 \(A_{1}B_{1}C_{1} - ABC\) 是正三棱柱，所以四边形 \(B_1BCC_1\) 是矩形. 连结 \(B_1C\) 交 \(BC_1\) 于 \(E\)，则 \(B_1E = EC\). 连结 \(DE\). 在三角形 \(AB_1C\) 中，\(AD = DC\) 且 \(B_1E = EC\)，故 \(DE \parallel AB_1\). 又 \(AB_1 \not\subset\) 平面 \(DBC_1\)，\(DE \subset\) 平面 \(DBC_1\)，所以 \(AB_1 \parallel\) 平面 \(DBC_1\).
\item 作 \(DF \perp BC\)，垂足为 \(F\)，则 \(DF \perp\) 平面 \(B_1BCC_1\). 连结 \(EF\)，则 \(EF\) 是 \(ED\) 在平面 \(B_1BCC_1\) 上的射影. 由（I）知 \(AB_1 \parallel DE\)，结合 \(AB_1 \perp BC_1\) 得 \(DE \perp BC_1\)，从而 \(EF \perp BC_1\). 故 \(\angle DEF\) 是二面角 \(\alpha\) 的平面角. 设 \(AC = 1\)，则 \(DC = \frac{1}{2}\). 因为三角形 \(ABC\) 是正三角形，在直角三角形 \(DCF\) 中，\(DF = DC \cdot \sin C = \frac{\sqrt{3}}{4}\)，\(CF = DC \cdot \cos C = \frac{1}{4}\). 取 \(BC\) 中点 \(G\)，由 \(EB = EC\) 得 \(EG \perp BC\). 在直角三角形 \(BEF\) 中，\(EF^2 = BF \cdot GF\). 又 \(BF = BC - FC = \frac{3}{4}\)，\(GF = \frac{1}{4}\)，故 \(EF^2 = \frac{3}{4} \cdot \frac{1}{4}\)，即 \(EF = \frac{\sqrt{3}}{4}\). 于是 \(\tan \angle DEF = \frac{DF}{EF} = 1\)，即 \(\angle DEF = 45\degree\). 故二面角 \(\alpha\) 为 \(45\degree\).
\end{enumerate}
\end{solution}

\begin{problem}
已知直角坐标平面上一点 \(Q(2, 0)\) 和圆 \(C\)：\(x^2 + y^2 = 1\)，动点 \(M\) 到圆 \(C\) 的切线长等于圆 \(C\) 的半径与 \(|MQ|\) 的和. 求动点 \(M\) 的轨迹方程，说明它表示什么曲线，并画出草图.

\bitmapfigure[width=0.45\linewidth]{img/1994/old_subject_science/q24_fig1.png}
\end{problem}

\begin{solution}
如图，设 \(MN\) 切圆于 \(N\)，圆的半径 \(|ON| = 1\)，则 \(|OM|^2 = |MN|^2 + |ON|^2 = |MN|^2 + 1\). 依题意，动点 \(M\) 组成的集合为 \(P = \{M \mid |MN| = |MQ| + 1\} = \{M \mid \sqrt{|OM|^2 - 1} = |MQ| + 1\}\). 设点 \(M(x, y)\)，则 \(\sqrt{x^2 + y^2 - 1} = \sqrt{(x - 2)^2 + y^2} + 1\). 整理得 \(2x - 3 = \sqrt{(x - 2)^2 + y^2} \geqslant 0\)，即 \(3x^2 - y^2 - 8x + 5 = 0\)（\(x \geqslant \frac{3}{2}\)）. 经检验，坐标适合这个方程的点都属于集合 \(P\)，故这个方程为所求的轨迹方程. 所求方程可化为 \(\frac{\left(x - \frac{4}{3}\right)^2}{\frac{1}{9}} - \frac{y^2}{\frac{1}{3}} = 1\)（\(x \geqslant \frac{3}{2}\)）. 它所表示的曲线是以点 \(\left(\frac{4}{3}, 0\right)\) 为中心，实轴在 \(x\) 轴上的双曲线的右支，顶点坐标为 \(\left(\frac{5}{3}, 0\right)\)，草图如下.
\solutionfigure{\bitmapfigure[width=0.5\linewidth]{img/1994/old_subject_science/q24_solution_fig1.png}}
\end{solution}

\begin{problem}
设 \(\{a_n\}\) 是正数组成的数列，其前 \(n\) 项和为 \(S_n\)，并且对于所有的自然数 \(n\)，\(a_n\) 与 \(2\) 的等差中项等于 \(S_n\) 与 \(2\) 的等比中项.

\begin{enumerate}
\item 写出数列 \(\{a_n\}\) 的前 \(3\) 项；
\item 求数列 \(\{a_n\}\) 的通项公式（写出推证过程）.
\end{enumerate}
\end{problem}

\begin{solution}
\begin{enumerate}
\item 由题意 \(\frac{a_n + 2}{2} = \sqrt{2S_n}\)（\(n \in \N\)），整理得 \(S_n = \frac{1}{8}\left(a_n + 2\right)^2\). 当 \(n = 1\) 时，\(S_1 = a_1\)，故 \(\frac{a_1 + 2}{2} = \sqrt{2a_1}\)，解得 \(a_1 = 2\). 当 \(n = 2\) 时，\(S_2 = a_1 + a_2 = 2 + a_2\)，故 \(8(2 + a_2) = \left(a_2 + 2\right)^2\)，即 \(\left(a_2 - 2\right)^2 = 16\). 由 \(a_2 > 0\)，解得 \(a_2 = 6\). 当 \(n = 3\) 时，\(S_3 = 8 + a_3\)，故 \(8(8 + a_3) = \left(a_3 + 2\right)^2\)，即 \(\left(a_3 - 2\right)^2 = 64\). 由 \(a_3 > 0\)，解得 \(a_3 = 10\). 故该数列的前 \(3\) 项为 \(2\)，\(6\)，\(10\).
\item 由（I）猜想数列 \(\{a_n\}\) 有通项公式 \(a_n = 4n - 2\). 下面用数学归纳法证明数列 \(\{a_n\}\) 的通项公式是 \(a_n = 4n - 2\)（\(n \in \N\)）. 当 \(n = 1\) 时，因为 \(4 \times 1 - 2 = 2\)，结论成立. 假设 \(n = k\) 时结论成立，即有 \(a_k = 4k - 2\). 由题意 \(\frac{a_k + 2}{2} = \sqrt{2S_k}\)，将 \(a_k = 4k - 2\) 代入得 \(2k = \sqrt{2S_k}\)，解得 \(S_k = 2k^2\). 由题意又有 \(\frac{a_{k+1} + 2}{2} = \sqrt{2S_{k+1}}\)，且 \(S_{k+1} = S_k + a_{k+1}\). 将 \(S_k = 2k^2\) 代入得 \(\left(\frac{a_{k+1} + 2}{2}\right)^2 = 2\left(a_{k+1} + 2k^2\right)\)，整理得 \(a_{k+1}^2 - 4a_{k+1} + 4 - 16k^2 = 0\). 由 \(a_{k+1} > 0\)，解得 \(a_{k+1} = 2 + 4k = 4(k + 1) - 2\). 这就是说，当 \(n = k + 1\) 时结论成立. 根据数学归纳法，结论对所有的自然数 \(n\) 成立，即通项公式为 \(a_n = 4n - 2\).
\end{enumerate}
\end{solution}
